\documentclass[10pt,twocolumn]{article}

\usepackage[spanish,es-nodecimaldot]{babel}
\usepackage[T1]{fontenc}
\usepackage[utf8]{inputenc}
\usepackage{geometry}
\usepackage{url}
\usepackage{hyperref}

\geometry{top=1.9cm,bottom=2cm,left=1.7cm,right=1.7cm,columnsep=0.65cm}
\hypersetup{colorlinks=true, urlcolor=blue, citecolor=black, linkcolor=black}
\setlength{\parindent}{1em}
\setlength{\parskip}{0pt}
\renewcommand{\thesection}{\Roman{section}}
\renewcommand{\thesubsection}{\arabic{section}.\arabic{subsection}}
\makeatletter
\renewcommand{\section}{\@startsection{section}{1}{\z@}%
  {-2.5ex \@plus -1ex \@minus -.2ex}%
  {1.3ex \@plus .2ex}%
  {\normalfont\large\bfseries\centering}}
\makeatother

\title{\textbf{Detección de sesgo mediático mediante análisis de framing:\\
arquitectura multiagente para noticias colombianas}}
\author{Abel Albuez Sanchez, Kelly Joane Leon Torres, Juan Camilo Torres Peña, Jesús David Romero Melo\\
\small\textit{Maestría en Ingeniería de Sistemas y Computación, Pontificia Universidad Javeriana, Bogotá D.C., Colombia}\\
\small Procesamiento de Lenguaje Natural, 2026}
\date{}

\begin{document}

\maketitle
\begin{abstract}
\textit{Este trabajo aborda la detección de sesgo mediático como un problema
de framing, distinto del análisis de sentimiento, en la cobertura colombiana
de un mismo evento noticioso. Se propone una arquitectura multiagente que
verifica el evento y compara la selección léxica, el encuadre de actores y
citas, y el estilo o énfasis temático. El prototipo se prueba sobre un corpus
curado de artículos y deja explícitamente pendientes las métricas cuantitativas
del corpus final, sin reportar resultados que aún no han sido calculados.}
\end{abstract}

\noindent\textbf{Index Terms---} media bias, framing analysis, multi-agent systems,
Spanish NLP, Colombian news.

\section{IDEA DE PROYECTO}

\subsection{Área y problema}

El proyecto se ubica en el área de procesamiento de lenguaje natural aplicado
al análisis de medios de comunicación y busca detectar y comparar sesgo
mediático en la cobertura colombiana de un mismo evento noticioso. El objeto
de estudio es el \emph{framing}: la manera en que la selección de palabras,
la caracterización de los actores y el énfasis temático pueden orientar la
interpretación de un hecho, incluso cuando distintos medios informan sobre el
mismo acontecimiento.

El problema no se formula como análisis de sentimiento. Una cobertura puede
emplear un tono aparentemente neutro y, aun así, construir encuadres distintos
mediante elecciones léxicas, atribución de roles, selección de citas y
organización del contexto. El prototipo propone un pipeline multiagente que
verifica primero que los artículos se refieran al mismo evento y después
compara tres dimensiones: léxico/framing, actores/citas y estilo/énfasis.

\subsection{Justificación}

Comparar artículos sobre un mismo evento permite separar las diferencias de
redacción de las diferencias producidas por tratar hechos distintos. Esta
comparación es pertinente para el contexto colombiano, donde el corpus
histórico que se está construyendo permitirá estudiar la cobertura de prensa a
mayor escala y comparar períodos de gobierno. El prototipo actual sirve como
prueba del método sobre pocos artículos y conserva una etapa explícita de
verificación para no forzar el análisis de fuentes que no cubren el mismo hecho.

La arquitectura se inspira en el sesgo por selección léxica y etiquetado
(\emph{bias by word choice and labeling}, WCL) de \cite{hamborg2020}, en el
análisis de framing centrado en personas de \cite{hamborg2023} y en la
viabilidad de herramientas de análisis de cobertura basadas en modelos de
lenguaje mostrada por \cite{mediabiasdetector2025}.

\subsection{Preguntas de investigación}

\begin{enumerate}
  \item ¿Qué diferencias de selección léxica y etiquetado aparecen entre
  medios colombianos al cubrir un mismo evento noticioso?
  \item ¿Cómo cambia el encuadre de los actores según las personas o
  instituciones citadas, los roles atribuidos y las voces que reciben espacio?
  \item ¿Qué diferencias de énfasis temático y estilo pueden observarse entre
  las coberturas y cómo complementan el análisis léxico y de actores?
\end{enumerate}

\section{ESTADO DEL ARTE}

El sesgo mediático ha sido abordado en PLN principalmente como un problema
de \emph{framing}, distinto del análisis de sentimiento: no se trata del
tono positivo o negativo de un texto, sino de cómo la selección de
palabras y el encuadre de los actores involucrados en un hecho puede
inducir interpretaciones distintas en el lector ante el mismo evento.

\cite{hamborg2020} define el concepto de \emph{bias by word choice and
labeling} (WCL): el fenómeno por el cual distintos medios refieren al
mismo concepto o actor usando términos diferentes que activan
connotaciones distintas (por ejemplo, ``manifestantes'' frente a
``vándalos'' para el mismo grupo de personas). Este trabajo establece el
problema que aborda el presente proyecto y da nombre explícito a la
categoría de sesgo que se busca detectar.

\cite{hamborg2023} extiende este marco con un método de
\emph{person-oriented framing analysis}: en vez de analizar el texto de
forma agregada, el análisis se centra en cómo cada artículo retrata
específicamente a las personas involucradas en el evento cubierto —a
quién se cita directamente, a quién se atribuyen acciones, y con qué
términos. Este enfoque es la base metodológica directa del componente de
actores y citas de la arquitectura propuesta en este proyecto.

Más recientemente, \cite{mediabiasdetector2025} presentan
\emph{Media Bias Detector}, una herramienta que usa modelos de lenguaje
de gran escala para anotar en tiempo casi real miles de artículos de
noticias, permitiendo comparar cobertura entre medios a nivel de tema,
tono y encuadre. Los autores señalan que los enfoques tradicionales de
PLN, como el conteo de n-gramas, no logran capturar el contexto completo
de un artículo, lo que motiva un enfoque basado en LLMs similar al
adoptado en este proyecto.

Este proyecto se ubica en la intersección de estas tres líneas: adopta la
definición de sesgo por selección léxica de \cite{hamborg2020}, el
enfoque centrado en actores de \cite{hamborg2023}, y la viabilidad de un
pipeline basado en LLMs demostrada por \cite{mediabiasdetector2025},
aplicándolos específicamente al corpus de medios colombianos.

\section{DESCRIPCIÓN DEL CORPUS}

\subsection{Fuente y curaduría}

La fuente del prototipo es una curaduría manual por evento. Cada registro
representa un artículo de un medio y se agrupan los registros que cubren un
mismo hecho puntual mediante el campo \texttt{evento\_id}. El archivo de ejemplo,
\texttt{corpus\_ejemplo.json}, contiene un evento sobre el fallo de la Corte
Constitucional frente a una reforma tributaria y tres artículos de medios
ficticios. Las diferencias de framing fueron diseñadas intencionalmente para
probar el pipeline; por tanto, este archivo no debe confundirse con una
colección de URLs periodísticas verificadas ni con el corpus histórico final.

El corpus histórico de mayor escala se construirá por separado en
\texttt{news-retrieval/}. En esta Entrega 1 no se reportan cifras de ese corpus porque
aún no está completo.

\subsection{Características y ejemplo}

El ejemplo disponible contiene un evento, tres medios y fechas comprendidas
entre el 14 y el 15 de octubre de 2024. El esquema de almacenamiento incluye
\texttt{evento\_id}, \texttt{medio}, \texttt{fecha}, \texttt{titular}, \texttt{texto} y \texttt{url}; el orquestador usa
\texttt{evento\_id} para agrupar artículos antes de ejecutar la verificación. Un
fragmento real del registro de \emph{El Diario Nacional} es:

\begin{quote}
La Corte Constitucional declaró exequible la reforma tributaria impulsada por
el Ejecutivo, en un fallo calificado por analistas como un espaldarazo a la
política social del Gobierno. La ministra de Hacienda celebró la decisión como
``un triunfo de la justicia social'' y aseguró que los recursos recaudados
financiarán programas de salud y educación para los sectores más vulnerables.
\end{quote}

Este fragmento ilustra por qué el corpus es útil para una prueba de framing:
la expresión ``espaldarazo'' y la cita sobre ``justicia social'' construyen un
encuadre diferente del que aparece, por ejemplo, en el titular que denomina
la reforma un ``impuestazo''.

\subsection{Limitaciones del corpus actual}

El corpus curado es pequeño y didáctico: reúne una sola noticia-evento y tres
fuentes ficticias. Además, la búsqueda web implementada para el modo real
recupera artículos sobre un tema, no garantiza que sean del mismo evento y
por eso siempre pasa por el agente verificador. Estas condiciones limitan la
posibilidad de generalizar los resultados del prototipo al conjunto de medios
colombianos.

\section{TEXTOMETRÍA BÁSICA}

\subsection{Limpieza y representación}

La limpieza se aplicará al campo \texttt{texto} después de conservar una copia del
registro original. El procedimiento contempla normalizar espacios y caracteres
problemáticos, segmentar el texto en oraciones y palabras, y conservar la
información necesaria para distinguir titulares, cuerpo, fechas y medios. La
normalización no debe eliminar las expresiones con carga de framing, por lo
que los términos originales se conservarán para el análisis cualitativo y
comparativo.

\subsection{Medidas textuales y frecuencias}

Las medidas básicas previstas incluyen número de palabras, número de
oraciones, longitud media de oración, diversidad léxica y frecuencias de
palabras y n-gramas. También se compararán las frecuencias por medio y por
evento, prestando atención a los términos que nombran actores, acciones y
conceptos centrales. Estas medidas describirán el corpus, pero no sustituirán
el análisis contextual de framing: un término frecuente no es necesariamente
un indicador de sesgo sin examinar el concepto al que refiere y su uso en el
artículo.

El corpus final todavía no está disponible con el tamaño suficiente para
reportar resultados cuantitativos. % TODO: completar con cifras reales del corpus final

\subsection{Visualización}

La visualización prevista incluye distribuciones de longitud de artículos y
oraciones, tablas o gráficos de frecuencia de términos por medio y
comparaciones de n-gramas relevantes. Para evitar que las visualizaciones
reduzcan el problema a sentimiento, se organizarán alrededor de actores,
conceptos y expresiones de framing identificados en el pipeline. % TODO: completar con cifras reales del corpus final

En la implementación actual, el agente de estilo/énfasis no calcula estas
medidas: solicita al modelo una estimación cualitativa sobre extensión
relativa, titular alarmista y presencia de contexto. Esta decisión es
provisional y está documentada en el código; debe reemplazarse por métricas
cuantitativas con \texttt{textstat} o \texttt{textdescriptives}, reutilizando el pipeline del
Taller 1, antes de presentar resultados definitivos.

Esta entrega documenta la idea, el fundamento académico, el formato del
corpus y el plan de textometría del prototipo. Los resultados numéricos de
frecuencias, medidas textuales y visualizaciones quedan pendientes hasta
contar con el corpus final y ejecutar el procesamiento cuantitativo
correspondiente. % TODO: completar con cifras reales del corpus final

\bibliographystyle{ieeetr}
\bibliography{referencias}

\end{document}
